\section{Measuring Similarity}
\label{sec:comparison}
\subsection{A Distance Between Problems}
The image models see five snapshots of each solution, at time indices
\[
 \mathcal I=\{0,25,50,74,99\},\qquad Q=5.
\]
We compare problems on these same snapshots. Let $X_i^{(r)}\in\mathbb R^{Q\times C\times H\times W}$ hold the five snapshots of sample $r$ of problem $i$, with $C=2$ fields, $H=W=64$, and $N=100$ samples. The reference distance is
\begin{equation}
 \boxed{D_{ij}=\left[
 \frac{1}{NQCHW}\sum_{r=1}^{N}
 \left\|X_i^{(r)}-X_j^{(r)}\right\|_F^2
 \right]^{1/2}.}
\label{eq:rms}
\end{equation}
This is the RMS difference between the two solutions, averaged over samples that share their initial data. A small value means that the two problems evolve similarly for these ICs, this time horizon, and these field scales. We also report a version that uses all 100 saved states. Five snapshots give a discrete average; we do not claim that they represent the full continuous-time behavior.

\begin{proposition}[Distance on observed behavior]
Stack the scaled snapshots of problem $i$ into one vector $F_i$. Then $D_{ij}=\|F_i-F_j\|_2$, so $D$ is symmetric, nonnegative, and satisfies the triangle inequality. It is a metric on distinct recorded solutions and a pseudometric on problems.
\label{prop:distance}
\end{proposition}
The second statement is a reminder that two equations could agree on these samples and times but differ elsewhere. The proof and the exact scaling are in \cref{app:trajectory-distance}.

\paragraph{Which pairs we compare.}
Our main comparisons hold the IC family fixed. Each family gives an $18\times18$ matrix covering the six equations and three BCs, all started from the same initial arrays. The full $54\times54$ matrix also compares problems from different IC families, so it mixes the effect of the equation with the effect of the initial data. We treat it as a secondary result.

\subsection{From Distances to Dot Products}
Embedding models report similarity as a dot product, so it helps to write the reference the same way. With the kernel $K_{ij}=F_i^\top F_j$,
\begin{equation}
 D_{ij}^2=K_{ii}+K_{jj}-2K_{ij}.
\label{eq:raw-kernel}
\end{equation}
This links the distance to an inner product while keeping the size of each field.

For an embedding model, let $z_i^{(r)}$ be the unit-length embedding of sample $r$ of problem $i$. We average the dot products over samples:
\begin{equation}
 S_{ij}^{\mathrm{emb}}=
 \frac1N\sum_{r=1}^{N}z_i^{(r)\top}z_j^{(r)}.
\label{eq:embedding-score}
\end{equation}
The matching embedding distance satisfies
\begin{equation}
 (d_{ij}^{\mathrm{emb}})^2
 =\frac1N\sum_r\|z_i^{(r)}-z_j^{(r)}\|_2^2
 =2-2S_{ij}^{\mathrm{emb}},
\label{eq:chord}
\end{equation}
so a larger dot product always means a smaller distance. We only ask whether $S^{\mathrm{emb}}$ ranks pairs in the same order as $-D$; the actual values need not match.

We also use a \emph{physical cosine}: normalize each simulated solution to unit length and then average the paired dot products, exactly as in \cref{eq:embedding-score}. This score ignores how large a response is and keeps only its shape. It can rank pairs differently from $D$, so we report it separately (\cref{app:physical-kernel}).

\subsection{Calibration With Exact Solutions}
For diffusion and linear exchange, each Laplacian mode evolves independently:
\begin{equation}
 \widehat\U_\ell(t)
 =\exp\!\big[t(\lambda_\ell\D+R)\big]
 \widehat\U_\ell(0),
\label{eq:linear-reference}
\end{equation}
with $R=0$ or $R=\kappa\left(\begin{smallmatrix}-1&1\\1&-1\end{smallmatrix}\right)$ \citep{higham2008functions}. The eigenvalues and eigenvectors of our grid Laplacian are known for each BC, so this formula gives the grid solution with no time-stepping error, apart from rounding. Advection--diffusion can be solved exactly in a similar way (\cref{app:linear-reference}).

A continuous example makes the scale of $D$ concrete. For the heat equation with zero Dirichlet BC and $u_0=\sin(\pi x)\sin(\pi y)$, the solution is $u_\nu(t,x,y)=e^{-2\pi^2\nu t}\sin(\pi x)\sin(\pi y)$. Over the five observed times, the RMS difference between two diffusion coefficients is
\begin{equation}
 D(\nu,\nu')=\frac12
 \left[\frac15\sum_{n\in\mathcal I}
 \big(e^{-2\pi^2\nu t_n}-e^{-2\pi^2\nu't_n}\big)^2
 \right]^{1/2}.
\label{eq:heat-example}
\end{equation}
\Cref{tab:calibration} gives values from this formula. They are exact calculations, not simulations or embedding results, and this IC is not one of the benchmark families. We later use the same four coefficients to test whether text embeddings notice a change of a single number.

\begin{table}[t]
\centering
\caption{Exact heat-equation distances from \cref{eq:heat-example}, relative to $\nu=0.01$. An embedding that tracks behavior should rank the candidates in this order.}
\label{tab:calibration}
\begin{tabular}{ccc}
\toprule
Candidate $\nu'$ & Five-frame RMS & 100-frame RMS \\
\midrule
0.011 & 0.005100 & 0.004895 \\
0.020 & 0.047511 & 0.045899 \\
0.050 & 0.152800 & 0.150418 \\
\bottomrule
\end{tabular}
\end{table}

\subsection{When Can We Trust a Ranking?}
Simulated distances contain numerical error, and a small error could swap two nearly equal distances. The following bound says when that cannot happen.
\begin{proposition}[A sufficient numerical margin]
Suppose $\|F_i-\widetilde F_i\|_2\le\epsilon_i$, where $\widetilde F_i$ is the computed solution. Then
\begin{equation}
 |D_{ij}-\widetilde D_{ij}|\le\epsilon_i+\epsilon_j.
\end{equation}
If $\widetilde D_{ik}-\widetilde D_{ij}>2\epsilon_i+\epsilon_j+\epsilon_k$, then the exact distances satisfy $D_{ij}<D_{ik}$.
\label{prop:margin}
\end{proposition}
The proof is the reverse triangle inequality. The bound is only as good as the error estimates $\epsilon_i$. For the linear equations we compute $\epsilon_i$ against exact solutions, so a ranking that passes the test is proved, up to rounding. For the nonlinear equations we estimate $\epsilon_i$ by refining the time step; this is strong evidence but not a proof.
